\documentclass[10pt,a4paper]{amsart}
\usepackage[margin=19mm]{geometry}
\usepackage{amsmath,amssymb,mathtools}
\usepackage[scr=rsfs,cal=euler]{mathalfa}
\usepackage{xfrac}
\usepackage[unicode,hidelinks,bookmarks=false]{hyperref}
\newcommand{\ie}{i.e.\ }
\newcommand{\cf}{cf.\ }
\newcommand{\resp}{respectively\ }
\newcommand{\Subsection}{\subsection}
\providecommand{\nobreakdash}{\nobreak\hbox{-}}
\setlength{\emergencystretch}{2em}
\multlinegap0pt
\usepackage{bm}
\usepackage{bbm}
\usepackage{dsfont}
\usepackage{xspace}
\usepackage{cancel}
\usepackage{mathtools}
\usepackage{amsthm}
\makeatletter
\@ifundefined{lemma}{\newtheorem{lemma}{Lemma}}{}
\@ifundefined{prop}{\newtheorem{prop}{Proposition}}{}
\makeatother
\newcommand{\Log}{\operatorname{Log}}
\newcommand{\SLERP}{\operatorname{SLERP}}
\newcommand{\bR}{\mathbb{R}}
\newcommand{\bS}{\mathbb{S}}
\title[Local Consistency and Higher-Order Structure of Spherical In]{Local Consistency and Higher-Order Structure of Spherical Interpolation}
\author{Shingyu Leung}
\date{Source: arXiv:2606.12802v1 (CC BY 4.0)}
\begin{document}
\maketitle

\noindent\emph{Source and licence.} Local Consistency and Higher-Order Structure of Spherical Interpolation. Version arXiv:2606.12802v1 (CC BY 4.0), distributed under the Creative Commons Attribution 4.0 International licence, as recorded in the arXiv licence metadata for this exact version and shown on the abstract page https://arxiv.org/abs/2606.12802v1. Full rights notice: licenses/arxiv-2606.12802-CC-BY-4.0.txt.\par\smallskip

\noindent\textit{Editorial scope.} Counted unit: the Prop whose source label is \texttt{prop:sider2\_refined} in the article (source0.tex lines 885--910, source label \texttt{prop:sider2\_refined}), delivered together with the article's own complete original proof of it (source0.tex lines 912--1102, one \texttt{proof} environment of 191 lines). No mathematics is written, repaired or re-derived: statement and proof are the article's own text line for line. Required context retained uncounted: eq:sider2 (lines 331--340); eq:curve\_taylor (lines 488--497); lem:cubic\_slerp (lines 523--556). Every definition, display and label of those lines is retained unchanged. Omitted from this package and not redistributed: 1--330, 341--487, 498--522, 557--884, 911--911, 1103--2614. Nothing inside a retained excerpt is omitted or altered: every retained body is delivered exactly as the article's lines are written, so no omission receipt of any kind accompanies this package. Citations: the retained excerpts carry no citation command at all, so no bibliography entry of the article is transcribed and no reference is redistributed. Reference adaptation: none was needed and none was made; every reference inside the delivered unit resolves inside it, so no unresolved reference or citation is produced. Printed slips kept literal: no printed word, symbol, hypothesis or number of the counted statement or of its proof was altered. Layout and class: the article's own class is \texttt{article} with packages including grffile, latexsym, amssymb, enumerate, amsmath, epsfig, amsthm, geometry, setspace, color, graphics, graphicx, algorithm2e, empheq; the wrapper is the lane's pinned \texttt{amsart} editorial wrapper at 10pt on A4, which restates the theorem-like environments the retained text needs and adds the article's own macro definitions that the retained text needs (\texttt{\textbackslash Log}, \texttt{\textbackslash SLERP}, \texttt{\textbackslash bR}, \texttt{\textbackslash bS}), with extra wrapper packages (bm, bbm, dsfont, xspace, cancel, mathtools). The delivered numbering is the unit's own. Assets: no retained span contains a figure environment, a graphics-inclusion command, a PDF-page inclusion, a table or a TikZ picture; no image file is copied into this package, the package declares no asset dependency and no image file is redistributed with it. Licence: version arXiv:2606.12802v1 of this article is distributed under the Creative Commons Attribution 4.0 International licence, as recorded in the arXiv licence metadata for this exact version and shown on the abstract page https://arxiv.org/abs/2606.12802v1. A full rights notice is delivered at licenses/arxiv-2606.12802-CC-BY-4.0.txt. Characters: every retained excerpt is pure ASCII, so the delivered engine, XeLaTeX with its default Latin Modern text font, reports no missing character.
\begin{equation}
\label{eq:sider2}
P_{012}(s)
=
\SLERP\left(
        \SLERP(p_0,d_a,\tau),
        \SLERP(d_b,p_2,\tau),
        \tau
       \right).
\end{equation}

\begin{equation}
\label{eq:curve_taylor}
x_i
=
(i-\theta)h V
+\frac{(i-\theta)^2h^2}{2}A
+\frac{(i-\theta)^3h^3}{6}B
+\frac{(i-\theta)^4h^4}{24}D
+O(h^5),
\end{equation}

\begin{lemma}[Cubic SLERP expansion in normal coordinates]
\label{lem:cubic_slerp}
Let \(y\in\bS^2\), and let \(A,B\) be sufficiently close to \(y\).  Put
\[
        U=\Log_y(A),\qquad W=\Log_y(B).
\]
For \(\lambda\) in a fixed bounded interval,
\begin{equation}
\label{eq:cubic_slerp}
\Log_y\!\left(\SLERP(A,B,\lambda)\right)
=
(1-\lambda)U+\lambda W+\mathcal C_\lambda(U,W)+O(\rho^5),
\end{equation}
where
\[
        \rho=|U|+|W|,
        \qquad
        Z=(1-\lambda)U+\lambda W,
\]
and
\begin{align}
\label{eq:cubic_term}
\mathcal C_\lambda(U,W)
= \frac{1}{6}\Big\{&
|Z|^2Z
+(1-\lambda)\lambda(2-\lambda)|U-W|^2U
+\lambda(1-\lambda^2)|U-W|^2W     \nonumber\\
&-(1-\lambda)|U|^2U
-\lambda |W|^2W
\Big\}.
\end{align}
In particular, the first non-affine correction is cubic in the endpoint
coordinates.
\end{lemma}

\begin{prop}[Refined local error of SIDER2]
\label{prop:sider2_refined}
Let \(K\subset\bR\) be a fixed bounded interval.  Let \(\gamma\in C^4\) on a
neighborhood of \(\{\theta h:\theta\in K\}\), and assume that all points
appearing in the SIDER2 construction remain in a common geodesically convex
ball of radius \(O(h)\).  Let \(P_{012}\) be the SIDER2 interpolant
\eqref{eq:sider2}.  For \(s=\theta h\), uniformly for \(\theta\in K\),
\begin{equation}
\label{eq:sider2_refined}
\Log_{\gamma(s)} P_{012}(s)
=
-\frac{h^3}{6}\,
\theta(\theta-1)(\theta-2)\, B
+O(h^4),
\end{equation}
where \(B=X^{(3)}(0)\) is the third Taylor coefficient of the normal-coordinate
curve
\[
        X(\eta)=\Log_{\gamma(s)}\gamma(s+\eta)
\]
at \(\eta=0\), as in \eqref{eq:curve_taylor}.  Consequently, on any fixed
bounded interval of the normalized parameter \(\theta\),
\[
        d_{\bS^2}\!\left(\gamma(s),P_{012}(s)\right)=O(h^3).
\]
\end{prop}

\begin{proof}
Fix \(s=\theta h\), set \(y=\gamma(s)\), and write
\[
        x_i=\Log_y(p_i),\qquad i=0,1,2.
\]
By the normal-coordinate Taylor expansion \eqref{eq:curve_taylor},
\begin{equation}
\label{eq:sider2_data_expansion}
x_i
=
(i-\theta)hV
+
\frac{(i-\theta)^2h^2}{2}A
+
\frac{(i-\theta)^3h^3}{6}B
+
O(h^4),
\end{equation}
uniformly for \(\theta\in K\).  Here \(V,A,B\in T_y\bS^2\) are bounded
vectors depending on the derivatives of \(\gamma\) at \(s\).

We now express the SIDER2 construction in these normal coordinates.  The two
spherical control points are
\[
        d_a=\SLERP(p_2,p_1,2),
        \qquad
        d_b=\SLERP(p_0,p_1,2).
\]
Let
\[
        \delta_a=\Log_y(d_a),\qquad
        \delta_b=\Log_y(d_b).
\]
Applying the cubic SLERP expansion of Lemma~\ref{lem:cubic_slerp} gives
\[
        \delta_a=2x_1-x_2+\mathcal C_2(x_2,x_1)+O(h^4),
\]
and
\[
        \delta_b=2x_1-x_0+\mathcal C_2(x_0,x_1)+O(h^4).
\]
For the purpose of identifying the coefficient of \(h^3\), the cubic term
\(\mathcal C_2\) depends only on the \(O(h)\) parts of its two arguments.  By
\eqref{eq:sider2_data_expansion}, these leading parts are scalar multiples of
the same vector \(V\):
\[
        x_i=(i-\theta)hV+O(h^2).
\]
Hence the leading-order endpoints in each control-point SLERP lie on the same
geodesic through \(y\).  For collinear normal-coordinate arguments, SLERP
coincides with one-dimensional affine interpolation along that geodesic, and
the cubic correction in Lemma~\ref{lem:cubic_slerp} vanishes.  Therefore
\begin{equation}
\label{eq:control_expansion}
        \delta_a=2x_1-x_2+O(h^4),
        \qquad
        \delta_b=2x_1-x_0+O(h^4).
\end{equation}

Let
\[
        \tau=\frac{s}{2h}=\frac{\theta}{2}.
\]
The two inner SLERP points in SIDER2 are
\[
        A_\tau=\SLERP(p_0,d_a,\tau),
        \qquad
        B_\tau=\SLERP(d_b,p_2,\tau).
\]
Denote their normal coordinates by
\[
        a_\tau=\Log_y(A_\tau),\qquad
        b_\tau=\Log_y(B_\tau).
\]
Using Lemma~\ref{lem:cubic_slerp} again, together with
\eqref{eq:control_expansion}, we obtain
\[
        a_\tau
        =
        (1-\tau)x_0+\tau(2x_1-x_2)+O(h^4),
\]
and
\[
        b_\tau
        =
        (1-\tau)(2x_1-x_0)+\tau x_2+O(h^4).
\]
The same reasoning as above applies: the possible cubic SLERP correction is
computed from the \(O(h)\) components of the endpoints, and these components
are collinear multiples of \(V\).  Thus no \(h^3\) geometric correction is
produced by the inner SLERP operations.

The final SIDER2 value is
\[
        P_{012}(s)=\SLERP(A_\tau,B_\tau,\tau).
\]
Let
\[
        z_\tau=\Log_y P_{012}(s).
\]
A final application of Lemma~\ref{lem:cubic_slerp} yields
\[
        z_\tau
        =
        (1-\tau)a_\tau+\tau b_\tau+O(h^4),
\]
again because the cubic correction generated by the leading \(O(h)\) parts
vanishes in the collinear case.  Substituting the expressions for \(a_\tau\)
and \(b_\tau\), we find
\begin{align*}
        z_\tau
        &=
        (1-\tau)\bigl[(1-\tau)x_0+\tau(2x_1-x_2)\bigr]  \\
        &\qquad
        +\tau\bigl[(1-\tau)(2x_1-x_0)+\tau x_2\bigr]
        +O(h^4)                                                   \\
        &=
        (1-2\tau)x_0
        +
        4\tau(1-\tau)x_1
        +
        (2\tau^2-\tau)x_2
        +O(h^4).
\end{align*}
Since \(\tau=\theta/2\), this becomes
\begin{equation}
\label{eq:sider2_affine_quadratic}
        z_\tau
        =
        (1-\theta)x_0
        +
        \theta(2-\theta)x_1
        +
        \frac{\theta(\theta-1)}{2}x_2
        +O(h^4).
\end{equation}
The first three terms are precisely the value at \(s=\theta h\) of the
Euclidean quadratic interpolant through the three normal-coordinate data
\[
        x_0,\quad x_1,\quad x_2,
\]
where the exact normal-coordinate value at \(s\) is zero.

It remains to evaluate the leading interpolation error.  Substituting the
Taylor expansion \eqref{eq:sider2_data_expansion} into
\eqref{eq:sider2_affine_quadratic}, the contributions from the constant,
linear, and quadratic Taylor terms cancel, as they must for quadratic
interpolation.  The first nonzero contribution comes from the cubic term:
\[
\begin{aligned}
        z_\tau
        &=
        \frac{h^3}{6}
        \left[
        (1-\theta)(-\theta)^3
        +
        \theta(2-\theta)(1-\theta)^3
        +
        \frac{\theta(\theta-1)}{2}(2-\theta)^3
        \right]B
        +O(h^4).
\end{aligned}
\]
The scalar coefficient simplifies to
\[
        -\theta(\theta-1)(\theta-2).
\]
Therefore
\[
        z_\tau
        =
        -\frac{h^3}{6}
        \theta(\theta-1)(\theta-2)B
        +O(h^4).
\]
Since \(z_\tau=\Log_{\gamma(s)}P_{012}(s)\), this proves
\eqref{eq:sider2_refined}.

Finally, in a fixed normal neighborhood the geodesic distance from
\(\gamma(s)\) to \(P_{012}(s)\) is exactly the norm of the logarithm vector:
\[
        d_{\bS^2}(\gamma(s),P_{012}(s))
        =
        \left|\Log_{\gamma(s)}P_{012}(s)\right|.
\]
The expansion above therefore implies
\[
        d_{\bS^2}(\gamma(s),P_{012}(s))=O(h^3),
\]
uniformly for \(\theta\in K\).
\end{proof}

\end{document}
